\documentclass[12pt]{article}
\usepackage[UTF8]{ctex}
\usepackage{amsmath,amssymb}
\usepackage{geometry}
\usepackage{graphicx}
\usepackage{hyperref}
\usepackage{booktabs}
\usepackage{array}
\geometry{margin=1in}

\newcommand{\E}{\mathbb{E}}

\begin{document}

\title{信息技术基础（2026 年春季）\\ 作业 \#3 \\ 线性神经网络、MDP 与启发式搜索}
\author{}
\date{}
\maketitle

%==============================================================
\section*{问题 1：线性激活神经网络（30 分）}

每个神经元采用线性激活函数
\[
g(x)\;=\;c + b\cdot\frac{1}{n}\sum_{i=1}^{n} W_i\,x_i,
\]
其中 $b,c$ 为固定实数，$n$ 为该神经元的入链数。
约定输入 $x = (x_0,x_1,\dots,x_n)$ 中 $x_0\equiv 1$ 充当偏置端（与 $W_0$ 配对），不影响题目形式。

%--------------------------------------------------------------
\subsection*{(1) 单个神经元：误差函数与梯度更新规则}

\paragraph{预测输出.}
\[
\hat{y}\;=\;g(\mathbf{x})\;=\;c+\frac{b}{n}\sum_{i=1}^{n}W_i\,x_i.
\]

\paragraph{平方误差函数.}
\[
\boxed{\;
E(\mathbf{W})\;=\;\tfrac{1}{2}\bigl(y-\hat{y}\bigr)^{2}.
\;}
\]
（系数 $\tfrac12$ 仅为求导方便，不改变极值点。）

\paragraph{梯度.}
由于 $\hat{y}$ 关于 $W_k$ 是线性的，
\[
\frac{\partial \hat{y}}{\partial W_k}\;=\;\frac{b}{n}\,x_k,
\qquad
\frac{\partial E}{\partial W_k}\;=\;-(y-\hat{y})\cdot\frac{b}{n}\,x_k.
\]

\paragraph{梯度下降权重更新规则.}
学习率 $\eta>0$，
\[
\boxed{\;
W_k \;\leftarrow\; W_k - \eta\,\frac{\partial E}{\partial W_k}
\;=\; W_k + \eta\,\frac{b}{n}\,(y-\hat{y})\,x_k,\qquad k=1,\dots,n.
\;}
\]
对 $W_0$（偏置权重）取 $x_0=1$，更新形式相同。

%--------------------------------------------------------------
\subsection*{(2) 一层隐藏层网络等价于无隐藏层网络}

\paragraph{隐藏层第 $j$ 个单元（入度 $n$）.}
\[
h_j \;=\; c + \frac{b}{n}\sum_{k=1}^{n}w_{k,j}\,x_k,\qquad j=1,\dots,m.
\]

\paragraph{输出单元（入度 $m$）.}
\[
o \;=\; c + \frac{b}{m}\sum_{j=1}^{m}W_j\,h_j.
\]

\paragraph{展开为关于 $w_{k,j},W_j,x$ 的表达式.}
代入 $h_j$ 并利用线性性：
\[
\begin{aligned}
o &= c + \frac{b}{m}\sum_{j=1}^{m}W_j\!\left(c + \frac{b}{n}\sum_{k=1}^{n}w_{k,j}\,x_k\right)\\[4pt]
  &= \underbrace{\,c + \frac{b\,c}{m}\sum_{j=1}^{m}W_j\,}_{=:A,\;\text{与 }x\text{ 无关}}
   \;+\;
     \underbrace{\,\frac{b^{\,2}}{m\,n}\sum_{k=1}^{n}\!\left(\sum_{j=1}^{m}W_j\,w_{k,j}\right)\!x_k\,}_{x\text{ 的线性项}}.
\end{aligned}
\]
记
\[
A\;=\;c+\frac{b\,c}{m}\sum_{j=1}^{m}W_j,
\qquad
B_k\;=\;\frac{b^{\,2}}{m\,n}\sum_{j=1}^{m}W_j\,w_{k,j},
\]
则
\[
\boxed{\;
o \;=\; A \;+\; \sum_{k=1}^{n} B_k\,x_k.
\;}
\]

\paragraph{结论：仿射函数可由单层线性网络精确表示.}
$o$ 是输入 $x$ 的仿射函数。考虑只含一个输出神经元、无隐藏层的网络
\[
o' \;=\; c + \frac{b}{n}\!\left(W'_0 + \sum_{k=1}^{n}W'_k\,x_k\right),
\]
（$x_0\equiv 1$ 为偏置端）。设 $b\neq 0$，取
\[
W'_k\;=\;\frac{b}{m}\sum_{j=1}^{m}W_j\,w_{k,j}\quad(k=1,\dots,n),
\qquad
W'_0\;=\;\frac{n}{b}(A-c)\;=\;\frac{n\,c}{m}\sum_{j=1}^{m}W_j.
\]
则代入直接验证：
\begin{itemize}
  \item 线性项：$\displaystyle\frac{b}{n}W'_k=\frac{b^{2}}{m\,n}\sum_j W_j w_{k,j}=B_k$；
  \item 常数项：$\displaystyle c+\frac{b}{n}W'_0=c+(A-c)=A$。
\end{itemize}
因此 $o'\equiv o$ 对一切输入成立。
若 $b=0$，则 $o\equiv c$ 与 $x$ 无关，单层网络（任意权重）恒输出 $c$，结论平凡成立。
$\blacksquare$

\paragraph{直觉.}
仿射映射的复合仍为仿射映射；线性激活神经元的任意层数堆叠都不会扩大假设空间——这正是深度网络必须使用非线性激活函数的根本原因。

\newpage
%==============================================================
\section*{问题 2：马尔可夫决策过程（30 分）}

约定 $T(s,a,s')=\Pr[s_{t+1}=s'\mid s_t=s,a_t=a]$；$R(s,a,s')$ 为转移即时奖励；折扣因子 $\gamma\in(0,1]$。

%--------------------------------------------------------------
\subsection*{(1) 总折扣奖励 $V$}

题目以 $t=1$ 为序列起点，但累计要求「$t=0\to+\infty$」，因此第一笔奖励的折扣指数取 $0$：
\[
\boxed{\;
V \;=\; \sum_{t=0}^{+\infty}\gamma^{t}\,R(s_{t+1},a_{t+1},s_{t+2})
   \;=\; \sum_{t=1}^{+\infty}\gamma^{t-1}\,R(s_{t},a_{t},s_{t+1}).
\;}
\]
展开形式：
$V=R(s_1,a_1,s_2)+\gamma R(s_2,a_2,s_3)+\gamma^{2} R(s_3,a_3,s_4)+\cdots$。

%--------------------------------------------------------------
\subsection*{(2) 值函数 $V^{\pi}(s_0)$}

策略 $\pi$ 给定下，从 $s_0$ 出发的期望累计折扣奖励：
\[
\boxed{\;
V^{\pi}(s_0)
\;=\; \E_{\pi}\!\left[\,
\sum_{t=0}^{\infty}\gamma^{t}\,R(s_{t+1},a_{t+1},s_{t+2})\;\Big|\;s_1=s_0
\,\right].
\;}
\]
期望对策略 $\pi$ 选择的（可能随机的）动作以及由 $T$ 给出的状态转移联合取得。

\paragraph{等价递归形式（Bellman 期望方程）.}
\[
V^{\pi}(s)\;=\;\sum_{a\in A}\pi(a\mid s)\sum_{s'\in S} T(s,a,s')\!\left[R(s,a,s')+\gamma\,V^{\pi}(s')\right].
\]

%--------------------------------------------------------------
\subsection*{(3) 值迭代（Value Iteration）求解 MDP}

\paragraph{目标.} 求最优值函数 $V^{*}$ 满足 Bellman 最优方程
\[
V^{*}(s)\;=\;\max_{a\in A}\sum_{s'\in S} T(s,a,s')\!\left[R(s,a,s')+\gamma\,V^{*}(s')\right].
\]

\paragraph{算法步骤.}
\begin{enumerate}
\item \textbf{初始化}：对所有 $s\in S$，置 $V_{0}(s)=0$（也可任意有限初值）。
\item \textbf{迭代}（Bellman 最优备份）：对 $k=0,1,2,\dots$ 与每个 $s\in S$，
\[
V_{k+1}(s)\;\leftarrow\;\max_{a\in A}\sum_{s'\in S} T(s,a,s')\!\left[R(s,a,s')+\gamma\,V_{k}(s')\right].
\]
\item \textbf{收敛判据}：当 $\bigl\|V_{k+1}-V_{k}\bigr\|_{\infty}<\varepsilon$（预设容差）时停止，记 $V^{*}\approx V_{K}$。
\item \textbf{导出最优策略}：
\[
\pi^{*}(s)\;=\;\arg\max_{a\in A}\sum_{s'\in S} T(s,a,s')\!\left[R(s,a,s')+\gamma\,V^{*}(s')\right].
\]
\end{enumerate}

\paragraph{收敛性.}
当 $\gamma<1$，Bellman 最优算子在最大模 $\|\cdot\|_{\infty}$ 下是 $\gamma$-收缩映射；由 Banach 不动点定理得 $V_{k}\to V^{*}$，且
\[
\|V_{k}-V^{*}\|_{\infty}\;\le\;\gamma^{k}\,\|V_{0}-V^{*}\|_{\infty},
\]
即指数（线性）速率收敛。$\gamma=1$ 时需额外条件（例如终止状态）保证良定义。

\newpage
%==============================================================
\section*{问题 3：启发式搜索（40 分）}

\textbf{可采纳启发式}：$h(s)\le h^{*}(s)$ 对所有状态 $s$ 成立，其中 $h^{*}(s)$ 为从 $s$ 到目标的最低真实代价（永不高估）。
设 $h_{1},h_{2}$ 均可采纳。

%--------------------------------------------------------------
\subsection*{(1) 哪些组合保证可采纳？（25 分）}

\begin{table}[h]
\centering
\renewcommand{\arraystretch}{1.45}
\begin{tabular}{>{\centering\arraybackslash}m{0.21\linewidth}c m{0.55\linewidth}}
\toprule
\textbf{组合} & \textbf{是否保证} & \multicolumn{1}{c}{\textbf{论证}}\\
\midrule
$h_1+h_2$            & \ding{55} & 反例：$h_1=h_2=h^{*}\Rightarrow h_1+h_2=2h^{*}>h^{*}$（$h^{*}>0$ 时高估）。\\
$h_1\times h_2$      & \ding{55} & 反例：$h_1=h_2=h^{*}=5\Rightarrow h_1 h_2=25>5$。\\
$\max(h_1,h_2)$      & \ding{51} & $\max(h_1,h_2)\le\max(h^{*},h^{*})=h^{*}$。\\
$\min(h_1,h_2)$      & \ding{51} & $\min(h_1,h_2)\le h_1\le h^{*}$。\\
$\alpha h_1+(1-\alpha)h_2$\newline $\alpha\in[0,1]$ & \ding{51} & 凸组合：$\le \alpha h^{*}+(1-\alpha)h^{*}=h^{*}$。\\
\bottomrule
\end{tabular}
\end{table}

\noindent
\textbf{答案}：保证可采纳的是 \textbf{(c)、(d)、(e)}。

\paragraph{补充说明.}
$\max(h_1,h_2)$ 是 A* 中常用的「pointwise 取大」技巧：在保持可采纳的前提下尽可能紧（dominating heuristic），通常能减少节点扩展数。
$h_1+h_2$ 不可采纳，但若两个启发式来自\emph{互不相交的子目标}（即不重复计算同一段代价），则它们的和仍可采纳——这就是 disjoint pattern database 的依据。

%--------------------------------------------------------------
\subsection*{(2) 判断题（15 分）}

\textbf{题目.} 考虑使用一致代价搜索（UCS）计算得到的最优路径。如果给每一步代价加上一个正常数，最优路径\emph{可能}发生改变。

\textbf{判断.} 正确（True）。

\paragraph{推理.}
UCS 选择路径总代价 $\displaystyle\mathrm{cost}(P)=\sum_{e\in P}\mathrm{cost}(e)$ 最小的路径 $P$。给每条边加 $c>0$ 后，
\[
\mathrm{cost}'(P)\;=\;\mathrm{cost}(P)+c\cdot|P|,
\]
其中 $|P|$ 为路径的边数。由于不同路径的 $|P|$ 通常不同，\textbf{额外增量 $c|P|$ 与路径长度耦合}：原本被「单边代价低但步数多」的路径所最优的情形，可能被「步数少」的路径反超。

\paragraph{显式反例.}
从起点 $s$ 到目标 $g$ 有两条候选路径：
\[
P_A\!:\;\text{$3$ 条边，每条代价 $4$}\;\Longrightarrow\;\mathrm{cost}(P_A)=12,
\qquad
P_B\!:\;\text{$1$ 条边，代价 $13$}\;\Longrightarrow\;\mathrm{cost}(P_B)=13.
\]
原始 UCS 选 $P_A$（$12<13$）。
加 $c=1$ 后：
\[
\mathrm{cost}'(P_A)=3\times 5=15,\qquad
\mathrm{cost}'(P_B)=14.
\]
新 UCS 选 $P_B$（$14<15$）。最优路径已改变。$\blacksquare$

\paragraph{直觉.}
UCS 的最优性对\emph{乘正常数}缩放是不变的（不改变排名），但对\emph{加正常数不是不变的}——加常数引入与路径长度耦合的额外项，等价于人为加入「按边数计费」的偏置，使最优解倾向更短（更少边）的路径。

\newpage
%==============================================================
\section*{答案速览}

\begin{table}[h]
\centering
\renewcommand{\arraystretch}{1.6}
\begin{tabular}{cl}
\toprule
\textbf{题号} & \textbf{关键结果}\\
\midrule
1(1) & $E(\mathbf{W})=\tfrac12(y-\hat{y})^{2}$\\
1(1) & $W_k\leftarrow W_k+\eta\,\dfrac{b}{n}(y-\hat{y})\,x_k$\\
1(2) & $o = A+\sum_{k=1}^{n}B_k\,x_k$，其中 $A=c+\dfrac{b\,c}{m}\sum_j W_j$，$B_k=\dfrac{b^{2}}{m\,n}\sum_j W_j\,w_{k,j}$\\
1(2) & 单层等价权重：$W'_k=\dfrac{b}{m}\sum_j W_j\,w_{k,j}$，$W'_0=\dfrac{n}{b}(A-c)$\\
2(1) & $V=\sum_{t=0}^{\infty}\gamma^{t}R(s_{t+1},a_{t+1},s_{t+2})$\\
2(2) & $V^{\pi}(s_0)=\E_{\pi}\!\left[\sum_{t=0}^{\infty}\gamma^{t}R(s_{t+1},a_{t+1},s_{t+2})\mid s_1=s_0\right]$\\
2(3) & $V_{k+1}(s)=\max_a\sum_{s'} T(s,a,s')[R(s,a,s')+\gamma V_k(s')]$；$\gamma$-收缩，速率 $\gamma^k$\\
3(1) & 保证可采纳：$\max(h_1,h_2)$、$\min(h_1,h_2)$、$\alpha h_1+(1-\alpha)h_2$，即 \textbf{(c)、(d)、(e)}\\
3(2) & \textbf{True}；反例：$P_A=3\times 4=12$，$P_B=13$；加 $c=1$ 后 $15>14$\\
\bottomrule
\end{tabular}
\end{table}

\end{document}
